\subsection{Information cutoffs require time within the day}
\label{sec:exact-information-time}

The assessment's civil date cannot determine whether a source was available
at the decision time. Evidence arriving at 00:00:01 on 1 April is later than a
cutoff at 00:00:00 on that date, although the two civil dates agree. The retained
bank assessment checks the full timestamp. The common requirement versions,
however, support ordered dates and rational numbers but have no timestamp type.
Passing a date comparison therefore cannot discharge that native timing duty.

Precision matters even when a timestamp parser is available. The retained native
schema accepts both \code{2026-04-01T00:00:00.0000001Z} and
\code{2026-04-01T00:00:00.0000002Z}. These instants differ by one tenth of a
microsecond, but its datetime-based conversion maps both to the same value.
That conversion gives the wrong ordering for a claim of exact timestamp
comparison. The retained bank examples themselves use whole-second timestamps;
their checked comparisons are unaffected. The defect concerns extending the
same conversion to all precision accepted by its schema.

The separate \code{ExactInstantProjection/1} addresses this boundary without
changing the original requirement evaluators. First, the existing binder selects
the timestamp strings from exact supplied record bytes. The projection then
converts each known value to rational seconds from midnight UTC on 1 January
of year 1, under an explicit proleptic Gregorian convention without leap seconds.
It counts whole days and seconds with integers, adds the exact decimal fraction,
and subtracts the supplied signed offset. No floating-point timestamp or
microsecond truncation enters this calculation. The resulting coordinate serves
ordering within the declared convention; physical elapsed time across leap
seconds is a separate question.

For example, 01:00:00 with offset \code{+01:00} and 00:00:00 with suffix
\code{Z} have equal coordinates. Their original strings and source identities
remain distinct. The two fractional instants above retain their different
rational coordinates, so the later one fails an available-by comparison against
the earlier cutoff. The existing dimensional requirement evaluator performs
that comparison. Each value has a seconds dimension and the same explicit
time-axis identity; the complete original context is preserved separately from
the time unit. A financial currency or scale in that context cannot become a
unit of time.

The accepted syntax requires a full date, hours, minutes and seconds, followed
by \code{Z} or an explicit signed hour-and-minute offset. An optional fraction
has at most 18 decimal digits. That precision bound is a declared reference
limit: more digits are refused, never rounded. Leap seconds, date-only values,
missing offsets, the unknown-offset spelling \code{-00:00}, named timezones and
UTC normalization outside years 1--9999 are unsupported. Invalid dates and
out-of-range clock or offset fields are invalid. A date-only disclosure needs
an explicit interval or availability policy before it could support an instant
comparison.

Unknown bindings keep their original reasons. Conversion preserves the selected
role, pointer, schema and raw-record identity together with the original string;
it cannot upgrade a supplied timestamp into an authenticated fact. Source-view
membership, the truth of availability and ingestion times, measurement dates
and the current view still require separate evidence. A capability to compare
times grants none of those admissions.

Independent integer and rational calculations cover 180 date/fraction points.
Six explicit offset and precision examples and 16 retained native timing
bindings agree with their stated comparators, and the associated text and typed
requirements agree. Malformed encodings, incompatible contexts and changed
bindings are refused. This optional conversion is adopted by the explicit
bank successor in Section~\ref{sec:exact-bank-profile}; other profiles retain
their own declared semantics. The counterexample and
exact conversion evidence are retained under
\path{data/accounting-master/exact-information-time-20261026/}.

\subsection{Exact timing must retain the whole evidence dependency}
\label{sec:exact-native-information}

Replacing timestamp arithmetic is useful only if the assessment still checks
every source on which its result depends. A current value may depend on an
earlier vintage whose evidence arrived too late. Checking only the current
record would admit that dependency even when its own timestamp comparison is
exact. The separately named \code{ExactBankInformation/1} therefore changes
only the retained auditor's information check. It walks the same complete
dependency closure and retains the original evidence-membership, ingestion
and measurement-date obligations.

For each evidence record in that closure, the checker requires its exact
reference to appear in the selected information view. Its available-at instant
must be no later than the decision cutoff, and its ingested-at instant no later
than the system cutoff. Its measurement date must be no later than the UTC
civil day containing the decision cutoff. That day is obtained from the exact
coordinate, including offset normalization. Unsupported or invalid timestamp
encodings are refused explicitly; there is no fallback to the lossy parser.

All 271 retained regulatory responses agree under the original and exact
auditors. Eleven additional information-check cases exercise later evidence,
ingestion, membership, measurement dates, equivalent offsets and precision.
Five outcomes intentionally change: later submicrosecond availability, later
submicrosecond ingestion and later superseded evidence are refused, as are
excess precision and the unknown-offset spelling \code{-00:00}. The superseded
case uses a native vintage record, which actually declares its predecessor.

These eleven examples exercise the information submethod; they are not eleven
new complete regulatory exchanges. A subsequent comparison checks ten complete
derived regulatory exchanges and four exact replays, including five intended
timing refusals. Its twelve adverse controls check the remaining request and
publication obligations. The original profile, packet and auditor remain
preserved. The separately identified profile below specifies conversion and
replay without changing those originals. Correct timing also leaves source
truth, signatures, state, budgets and accounting applicability to their existing
checks. The complete comparisons and derived examples are retained under
\path{data/accounting-master/exact-native-timing-20261026/}.

\subsection{Changing the timing rule creates a new bank profile}
\label{sec:exact-bank-profile}

A saved assessment records which rules were applied to which information.
Changing the time comparator must therefore leave the old assessment
interpretable. Consider evidence available one ten-millionth of a second after
the decision cutoff. The original parser can erase that difference and accept
the request; the exact comparator refuses it. Returning an old cached answer
under the new interpretation would conceal the change, even if the arithmetic
in the assessment were otherwise identical.

The documentary successor \code{BankAssessmentExactTime/1} gives that change an
explicit owner. It supports the same three operations: regulatory assessment,
conditional-loss query and policy description. Every evidence dependency must
meet the exact availability and ingestion comparisons described above.
Measurement, assessment and vintage/origin dates use the UTC civil day of the
exact decision instant. The parser retains the stated Gregorian calendar, known
offset and eighteen-digit precision limits; unsupported precision is refused.
The capital and loss methods retain their financial calculations. Only their
two decision-day conversions change.

The new profile binds the exact schema and semantic-description bytes. Eight
record kinds receive new schema identities because their typed references
depend on the profile: the profile itself, operation signature, request,
result, transport observation, query receipt, coverage domain and coverage
region. Unchanged financial records retain their original schemas and, when
their contents are unchanged, their original identities. A conversion creates
a new request in a distinct workspace and new exact submission bytes. It
publishes a new query receipt while preserving the old request, receipt and
query-index entry. Replacing a supplied evidence timestamp additionally changes
the identities of the records that depend on it.

Profile, schema, semantic-description and signature checks precede replay.
An old request consequently cannot acquire the successor's meaning merely
because it has a cached result. Within the successor, submitting the same
request bytes returns the same output, receipt, statuses and coverage, even
when the current view or accounting head has since advanced. Different bytes
under an already used request identity are refused. A new request still needs
the declared view, accounting head, execution environment and budget. A refusal
has no result, receipt or coverage and changes no world state. A successful or
partial query changes the immutable query records and index; it posts no
journal and changes no accounting balance or owner head.

The finite conversion is checked against all 33 retained positive publications:
25 regulatory assessments, seven conditional-loss queries and one policy
description. Expected outputs come from the saved publications with declared
reference substitutions. The comparisons preserve entire outputs, including
unknown amounts, undefined ratios, diagnostic codes, unresolved obligations
and coverage. They do not compare only a selected ratio. Each converted
publication coexists with its original query and receipt.

Thirty complete timing cases then exercise all three operations. Twenty-one
outcomes intentionally change relative to a comparator that retains the old
time methods: later fractional availability or ingestion, excessive precision
in evidence or cutoff times, unknown offsets and eighteenth-digit ordering
are refused. Equal times and equivalent offsets remain accepted. Forty-two
exact replays and 28 additional controls check identity conflicts, stale locks,
signatures, decoder descriptions, heads, views and publication effects. All 63
saved conversion packages reconstruct and pass their full response checks.

These examples establish a bounded conversion between two documentary
profiles. They share the retained financial methods and synthetic inputs; they
do not independently validate those methods against current banking rules or
authenticate a real disclosure's publication time. Concurrent service
execution, hardware-loss recovery, general profile migration and empirical
forecast accuracy require separate evidence. The original profiles and examples
remain preserved. The full-request comparisons are stored under
\path{data/accounting-master/exact-bank-request-conformance-20261026/};
the successor contract, converted exchanges and checks are under
\path{data/accounting-master/exact-bank-profile-adoption-20261026/}.


\subsection{Writing the complete method behind an event}
\label{sec:complete-method-authoring}

A command that selects a registered payment method does not yet explain how an
author writes that method. The payment needs a condition on available cash, a
decrease in the cash account, a decrease in outstanding principal, a financing
cash-flow explanation and a check that the resulting liability still agrees
with the loan record. Leaving any of those parts behind a method name would
leave a material part of the financial algorithm outside the authored text.

\code{CompanyMethodAuthoring/1} writes the complete retained
\code{CompanyBookMethod/1} definition. Its header identifies the method and
selects \code{CompanyBookTransition/1} or \code{/2}. The remaining declarations
supply context and account mappings, guarded preconditions with their input
bindings, monetary calculations and effects, ordered memory changes, guarded
postconditions and residual prerequisites. The original validators still
determine the meaning of each part. The text adds a way to author their inputs;
it does not introduce another accounting evaluator.

The transition version is consequential. Version 1 supports replacement of
existing scalar memory fields, such as outstanding principal. Version 2 also
checks typed insertion, allowing a credit sale to create an identified invoice
claim. Its compilation environment must therefore include the memory type
definition. Both versions require an explicitly supplied book, event and schema
set. Writing a method does not reconstruct those inputs from a company name or
from its disclosures.

Each declaration occupies one line and ends with a semicolon. Precondition and
postcondition sections contain the existing requirement-exchange text, with
empty event values and capabilities. Their selected content digests remain
checked. The monetary section declares inputs, calculations, account roles,
postings and gross flows. Expressions use references, constants, addition,
subtraction, multiplication and minimum. Types and dependencies are explicit;
the native monetary validator refuses an undeclared or forward reference,
an incomplete dependency declaration or an incompatible context. The memory
section names its schemas, binds its inputs and declares each ordered write.
Binding records, types and memory templates use strict JSON data tokens.

For example, the payment's outstanding-principal calculation reads the
existing principal and subtracts the supplied payment amount. A memory input
then names that calculation as the value to write back. The liability posting
uses the same payment amount with the book's signed-account convention, and
the cash posting and financing flow both record its negative. A postcondition
checks that liability plus outstanding principal is zero. The text therefore
exposes the relationship between the accounting effect and the continuing
contract state, rather than merely printing a balanced journal.

The checked full payment starts with cash 100 and liability $-30$, pays
$7/2$, and finishes with cash $193/2$ and liability $-53/2$. A separately
authored variant changes every monetary use of the payment to one half of that
amount. Its cash becomes $393/4$, liability $-113/4$, outstanding principal
$113/4$ and financing cash flow $-7/4$. The dimensionless multiplier retains
the payment's context, as required by the existing product rule. This
constructed variant tests whether text changes the algorithm itself; it does
not propose a rule for how much an actual borrower should pay.

Four complete payment and claim-creating sale definitions, using two company
charts, match their typed definitions, requests and full native results.
The sales increase receivables by 30, reduce inventory by 10 and insert the
identified claim with remaining amount 30; they create no cash flow. The
comparisons include balances, memory, journals and gross flow records. Twenty
syntax or environment refusal cases and twelve economic-outcome cases retain
the expected distinctions: an unknown amount is unresolved, a negative amount
is invalid, and false applicability is unsupported. A forged proposed
after-state is refused by the existing application step.

Compilation itself executes nothing. Successful syntax lowering is distinct
from compilation against an environment, and neither decision establishes
source or accounting admission. Exact UTF-8 source slices retain their own
hashes; whitespace and reordering of mapping declarations leave the semantic
identity unchanged. Ordered calculations and writes keep their order. The
bounded reference limits remain explicit, including one-megabyte text and
4096-node, 48-depth checks. The complete grammar, authored examples and results
are retained under
\path{data/accounting-master/complete-method-authoring-20261026/}.

This supplies complete authoring for the retained method definition. Broader
method algorithms, registry loading, company-state reconstruction and independent
implementation remain separate work. In particular, agreement with a shared
native engine is evidence of consistent compilation, not independent proof of
every accounting treatment implemented by that engine.

\subsection{From an authored method to a sequence of company statements}
\label{sec:authored-program-composition}

A credit sale creates a claim; its collection later needs that claim to exist.
The complete method text does not by itself determine when a method may be
called or which method version a command selects. The retained registry version
2 still routes its calls through transition version 1, which supports fixed
scalar replacement. Its version number therefore cannot be used to infer
support for claim insertion.

The separate \code{AuthoredMethodRegistry/1} records the transition version
explicitly for each entry, along with the exact authored source, its hash and
the complete method definition and hash. Lowering that source must reproduce
the same method and version. Event-input types, account classes, metadata
bindings and the direct-event-date clock remain explicit. A transition-v2
entry also supplies its memory type definition. The registry validates every
entry against the current book, including entries not selected for the call;
it then dispatches the selected event through the original transition engine.

This state dependence matters when composing commands. A collection method
that reads an identified claim cannot be treated as usable before the sale
creates that claim. \code{AuthoredCompanyProgram/1} therefore supplies a finite
population of explicitly hashed registries and lets each event name the one
it uses. The program first checks its complete command structure and source
identities. At each event, the selected registry is validated against the book
that actually exists at that point. Earlier commands never receive the benefit
of a future state merely to make validation succeed.

Consider the constructed lifecycle used in the preceding claim examples.
The sale increases receivables by 30 and creates the claim. A first receipt of
10 raises cash from 100 to 110 and leaves receivables and the claim at 20.
Due-date assessments change the declared claim observations without creating
cash. The annual report and period carry preserve that remaining claim while
moving the year's profit into retained earnings under the supplied closing
method. A receipt of 20 in the next year brings cash to 130 and clears both
receivables and the claim. Next-period operating cash flow is 20. Two company
charts reproduce this eleven-command sequence, including complete journals,
memory, reports and the period-carry result.

Report and carry commands call the existing native methods with explicit
cutoffs, mappings, equity movements, notes and per-leaf memory policies. The
complete requests and results agree with the retained lifecycle, not only its
cash totals. Method-input schemas remain distinct resources from the schemas
supplied to those report requests. Every successful step also preserves the
typed memory definition. Observing a book at a cutoff does not establish that
all economically relevant events have been supplied.

The comparison has an important representational qualification. The original
lifecycle builder specializes claim identity, customer and due date into its
generated methods, while retaining redundant event fields. The registry
requires exactly the declared inputs, so the derived comparison records every
removed undeclared field and the changed request identity. Complete states,
journals and flows still match the original lifecycle. Exact event-request
comparison is against direct execution of the explicitly derived request.
These methods do not yet supply arbitrary dynamic claim selection.

Failures preserve what happened earlier. If the first sale succeeds and the
next collection amount is unknown, the result retains the sale, its claim and
the unexecuted command suffix, with the collection unresolved. Negative amounts,
false applicability, repeated identities, early reporting cutoffs, incomplete
carry policies and invalid unselected entries retain their distinct refusals.
Seven such cases preserve the exact successful prefix. Forged complete results
fail recomputation. A separate authored edit halves the collection algorithm's
payment: after the sale and that receipt, cash is 105 and receivables and the
claim are 25. This intervention tests execution of the supplied algorithm.

The checked registry and program remain separate reference successors. They
preserve the original engines and their provisional status. The program's
64-command and eight-mebibyte trace limits are reference resource bounds,
not production capacity claims. Full textual authoring of registry resources,
general method specialization, broader financial methods and actual-company
source admission remain open. The comparisons are retained under
\path{data/accounting-master/authored-method-registry-20261026/} and
\path{data/accounting-master/authored-program-composition-20261026/}.

\subsection{A typed claim can still disagree with the book}
\label{sec:authored-state-invariants}

The preceding successful histories do not establish that every authored method
preserves the original lifecycle's economic controls. Suppose an otherwise
unchanged sale posts receivables of 30 but creates a claim whose remaining
amount is 31. Both values have the right monetary type, the journal balances,
and the memory record satisfies its schema. The original lifecycle nevertheless
rejects that state because the claims do not sum to receivables. Typed-book
validation alone does not express this cross-record relation.

A focused audit changes one claim field at a time in the authored sale and
then requests a report from the genuine period opening. Across two company
charts, sixteen cases expose eight outer-control gaps: claim totals, the sale
event's membership in the book, issue/due chronology, future issue dates,
negative assessment amounts, nonzero unassessed placeholders, future assessment
and an incorrect overdue snapshot. An empty customer string is already refused
by the memory schema. That covered structural rule must not be counted as a
missing economic invariant.

\code{BookStateInvariant/1} supplies the missing connection to the existing
collection evaluators. Its explicit bindings select the book's date and
receivable account, the identified rows of its claim map and its seen-event
identities. Each scalar's type and context must agree with the declared input,
including for an empty map. Values come from that exact supplied book; a caller
cannot substitute a Boolean assertion that its controls passed.

The unchanged collection evaluator checks unique claim identities, references
from each claim to its sale event and guarded conditions on every row. The
unchanged collection-total evaluator sums remaining amounts and compares that
total with receivables. An assessed claim must have an assessment date between
issue and the current date. Its overdue snapshot equals the assessed amount
when assessment occurs strictly after the due date, and zero otherwise.
An unassessed claim retains zero assessment and overdue placeholders. Remaining,
assessed and overdue amounts are nonnegative. These are declared controls of
this constructed lifecycle, not a new impairment model.

The separate \code{GuardedAuthoredCompanyProgram/1} selects the invariant by
its exact hash. It checks the opening, each current book and each proposed
after-state. Only a proposed state that satisfies the selected controls enters
the successful prefix. In the 31-versus-30 example, the proposed sale is
refused, the opening book remains current, and the sale and report remain in
the unexecuted suffix. A failed post-state check cannot quietly retain the
invalid claim or its otherwise balanced journal.

Both eleven-command histories retain their complete original results, and
all 24 opening and after-state books pass the original lifecycle control.
All sixteen changed-method counterexamples are refused. Wrong pointers,
contexts and field types also refuse. The original evaluators retain unresolved
results for partial or unknown collection inputs, while an unrelated unknown
does not erase a known contradiction. Unsupported capabilities and explicitly
unsupported consequences remain distinct from invalid states. Neither a
successful invariant nor a checked schema establishes source or accounting
admission.

Completeness here means all members of the supplied internal map. It does not
mean that an actual issuer disclosed every claim. Event-specific duties also
remain: selecting the right customer, preventing an inappropriate new event
and respecting sealed assessment windows require their own prerequisites.
The state-control repair therefore preserves the finite comparisons without
claiming whole-lifecycle substitutability. Counterexamples and repair evidence
are retained under
\path{data/accounting-master/authored-program-invariant-audit-20261026/} and
\path{data/accounting-master/authored-state-invariants-20261026/}.

\input{sections/authored-event-contract}



